\documentclass[12pt,letterpaper]{article}
\usepackage[utf8]{inputenc}
\usepackage[spanish]{babel}
\usepackage[margin=2.5cm]{geometry}
\usepackage{amsmath,amssymb}
\usepackage{tcolorbox}
\usepackage{enumitem}
\usepackage{hyperref}

\hypersetup{
    colorlinks=true,
    linkcolor=blue,
    urlcolor=blue
}

\tcbuselibrary{skins,breakable}

\newtcolorbox{promptbox}[1]{
    colback=blue!5!white,
    colframe=blue!75!black,
    fonttitle=\bfseries,
    title=#1,
    breakable,
    arc=2mm
}

\newtcolorbox{autobox}{
    colback=green!5!white,
    colframe=green!50!black,
    fonttitle=\bfseries,
    title={Lo que el Agente hará automáticamente tras recibir tu prompt},
    breakable,
    arc=2mm
}

\title{\textbf{Guía de Prompts y Flujo de Trabajo para Evaluación de Exámenes en Lote}\\[0.3em]
\large Física para Ciencias de la Salud (005-1713) --- Bioanálisis}
\author{Universidad de Oriente --- Núcleo de Sucre}
\date{\today}

\begin{document}

\maketitle

\section{Introducción}
Este documento contiene la guía estandarizada de \textit{prompts} e instrucciones para la evaluación automática y masiva de exámenes de la asignatura \textbf{Física para Ciencias de la Salud} mediante agentes de IA. Está diseñado para ser reutilizado en nuevos semestres y asignaturas afines.

\section{Plantillas de Prompts para Evaluación en Lote}

\begin{promptbox}{Opción 1: Prompt General para una Unidad Completa (Recomendado)}
Por favor evalúa en lote todos los exámenes contenidos en la carpeta \texttt{examenes/05/examen\_estudiantes}. Son [N] preguntas con un valor de [X] ptos cada una. Utiliza la rúbrica/solucionario de \texttt{examenes/05/solucionario\_examen\_unidad\_V.tex}. Al finalizar, compila los informes en PDF, limpia los archivos auxiliares de LaTeX y muéstrame la tabla consolidada de resultados.
\end{promptbox}

\vspace{0.5em}

\begin{promptbox}{Opción 2: Prompt con Criterios Específicos de Evaluación}
Vamos a evaluar en lote la carpeta \texttt{examenes/03/examen\_estudiantes}. Consta de [N] problemas con un valor de [X] ptos cada uno. Sigue estas instrucciones de corrección: ``\textit{Este examen evalúa la comprensión profunda y el análisis físico aplicado a fenómenos biológicos y clínicos. No basta con escribir fórmulas matemáticas; debe argumentar, analizar y justificar cada una de sus respuestas basándose en los principios físicos estudiados}''. Usa el modelo \texttt{gemini-3.5-flash} y compila los informes PDF limpiando los archivos temporales.
\end{promptbox}

\vspace{0.5em}

\begin{promptbox}{Opción 3: Prompt Mínimo / Rápido}
Evalúa en lote los exámenes de \texttt{examenes/02/examen\_estudiantes} usando el solucionario de esa unidad. Genera los informes PDF y limpia los archivos auxiliares.
\end{promptbox}

\section{Flujo Automático del Agente}

\begin{autobox}
Al recibir cualquiera de los \textit{prompts} anteriores, el agente ejecutará de forma autónoma el siguiente flujo multietapa:

\begin{enumerate}[leftmargin=*]
    \item \textbf{Construcción / Lectura de la Rúbrica:}
    Extrae los criterios de corrección, respuestas cuantitativas de referencia y deducciones físicas cualitativas requeridas a partir del solucionario máster (\texttt{.tex}) de la unidad.

    \item \textbf{Evaluación Multimodal con Conmutación Dinámica de Modelos:}
    Itera secuencialmente sobre los archivos PDF de la carpeta \texttt{examen\_estudiantes/}, convirtiendo las páginas manuscritas a imágenes para analizarlas con el modelo multimodal. Si un modelo agota su cuota diaria o de frecuencia (\texttt{429 RESOURCE\_EXHAUSTED}), el script conmuta automáticamente en cascada a un modelo secundario (\texttt{gemini-2.5-flash} $\rightarrow$ \texttt{gemini-3.5-flash} $\rightarrow$ \texttt{gemini-3.6-flash}) sin detener la ejecución.

    \item \textbf{Persistencia en JSON:}
    Almacena los puntajes desglosados por pregunta, observaciones pedagógicas, errores detectados y nivel de desempeño en un archivo JSON en \texttt{.tmp/evaluacion\_[Estudiante].json}.

    \item \textbf{Generación e Integración del Informe LaTeX:}
    Convierte el resultado JSON en un documento LaTeX individual estandarizado (\texttt{informe\_[Estudiante].tex}) ubicado en \texttt{informe\_examenes/}.

    \item \textbf{Compilación y Limpieza de Auxiliares:}
    Compila el informe a PDF usando \texttt{pdflatex} / \texttt{latexmk} y ejecuta el script determinista \texttt{clean\_latex.py} para eliminar archivos auxiliares (\texttt{.aux}, \texttt{.log}, \texttt{.out}, \texttt{.fls}, \texttt{.fdb\_latexmk}).

    \item \textbf{Presentación de Resultados Consolidados:}
    Muestra la tabla estadística final con los puntajes de cada estudiante, promedio del curso, porcentaje por niveles de desempeño y enlaces directos a cada informe PDF generado.
\end{enumerate}
\end{autobox}

\section{Estructura Recomendada de Archivos para Nuevos Semestres}

\begin{verbatim}
examenes/
+-- XX/                          # Número de Unidad (ej. 01, 02, ..., 05)
    |-- examen_estudiantes/      # PDFs de exámenes entregados por los estudiantes
    |-- solucionario_unidad_XX.tex# Solucionario máster en LaTeX
    |-- rubrica_unidad_XX.txt    # Rúbrica detallada (opcional)
    +-- informe_examenes/        # Informes PDF individuales producidos
\end{verbatim}

\end{document}
